% Recommended papers: Tools and datasets for mining libre software repositories
%			The promises and perils of mining git
%
%
%
%
%

% $Id: slides.tex 4228 2006-06-21 21:55:12Z jjamor $
%
%
% Compilar a .pdf con LaTeX (pdflatex)
% Es necesario instalar Beamer (paquete latex-beamer en Debian)
%

%
% Gráficos:
% Los gráficos pueden suministrarse en PNG, JPG, TIF, PDF, MPS
% Los EPS deben convertirse a PDF (usar epstopdf)
%

\documentclass{beamer}
\usetheme{Warsaw}
\usebackgroundtemplate{\includegraphics[width=\paperwidth]{format/libresoft-bg.png}}
\usepackage[spanish]{babel}
\usepackage[utf8]{inputenc}
\usepackage{graphics}
\usepackage{amssymb} % Simbolos matematicos
\usepackage{url}

%\definecolor{libresoftgreen}{RGB}{162,190,43}
%\definecolor{libresoftblue}{RGB}{0,98,143}

%\setbeamercolor{titlelike}{bg=libresoftgreen}

%% Metadatos del PDF.
\hypersetup{
  pdftitle={Tools to analyze libre software repositories},
  pdfauthor={Daniel Izquierdo},
  pdfcreator={GSyC/Libresoft},
  pdfproducer=PDFLaTeX,
  pdfsubject={nn},
}
%%


\AtBeginSection[]
{
  \begin{frame}<presentation>
    \frametitle{Index}
    \tableofcontents[current]
  \end{frame}
}


\begin{document}

\title{Tools to analyze Libre Software Repositories}
\subtitle{Master on Free Software}
\institute{\\dizquierdo@libresoft.es\\
GSyC/Libresoft}
\author{Daniel Izquierdo}
\date{\today}

\frame{
\maketitle
\begin{center}
\includegraphics[width=6cm]{format/gsyc-urjc}
\end{center}
}


% Si el titulo o el autor se quieren acortar para los pies de página
% se pueden redefinir aquí:
%\title{Titulo corto}
%\author{Autores abreviado}


%% LICENCIA DE REDISTRIBUCION DE LAS TRANSPAS
\frame{
~
\vspace{4cm}

\begin{flushright}
{\tiny
(cc) 2011 Daniel Izquierdo. \\
Some rights reserved. This document is distributed under the Creative \\
            Commons Attribution-ShareAlike 3.0 licence, available in \\
            http://creativecommons.org/licenses/by-sa/3.0/

%  Este documento (o uno muy similar) está disponible en \\
%  \url{http://gsyc.escet.urjc.es/~jjamor/}
}
\end{flushright}
}
%%

%%%%%%
%Transpas separadas por \begin{frame}
%%%%%%%%%%%%%%%%%%%%%%%%\end{frame}

\section{Introduction}

\begin{frame}
\frametitle{Data sources}
\begin{itemize}
\item Source code management system
\item Mailing lists
\item Bug tracking system
\item Source code
\end{itemize}
\end{frame}

%%%%%%%%%%%%%%%%%%%%%%%%%%%%%%%%%%%%%%%%%%%%%%%%%%%%%%%%%%%%%%

\begin{frame}
 \frametitle{Data sources}
 \begin{itemize}
 \item What type of tools are used?
 \end{itemize}
\end{frame}

%%%%%%%%%%%%%%%%%%%%%%%%%%%%%%%%%%%%%%%%%%%%%%%%%%%%%%%%%%%%%%


\section{Tools}


% to be done

\begin{frame}
 \frametitle{SCM analyzers}
 \begin{itemize}
  \item CVSAnalY: \url{http://git.libresoft.es/cvsanaly/}
  \item This tool retrieves information from SVN, CVS or Git and 
  creates a relational database.
 \end{itemize}
\end{frame}

%%%%%%%%%%%%%%%%%%%%%%%%%%%%%%%%%%%%%%%%%%%%%%%%%%%%%%%%%%%%%%

\begin{frame}
 \frametitle{Mailing lists analyzers}
 \begin{itemize}
 \item Mailing List Stats: \url{http://git.libresoft.es/mailingliststat/}
 \item This tool retrieves the headers and body of the e-mails and
 create a relational database.
 \end{itemize}
\end{frame}



\begin{frame}
 \frametitle{BTS analyzers}
 \begin{itemize}
 \item Bicho: \url{http://git.libresoft.es/bicho/}
 \item This tool retrieves information from the Bugzilla (KDE, Apache and GNOME),
 the SourceForge tracker and Jira. 
 \end{itemize}


\end{frame}



\begin{frame}
 \frametitle{Source Code analyzers}
 \begin{itemize}
  \item Basic metrics such as programming languages,
  number of SLOCs, number of files, functions, but independent from the programming language of the source code
   \begin{itemize}
      \item SLOCCount: \url{http://www.dwheeler.com/sloccount/}
      \item Cloc: \url{http://www.dwheeler.com/sloccount/}
      \item Ohcount: \url{http://www.dwheeler.com/sloccount/}
   \end{itemize}
 \end{itemize}
\end{frame}


\begin{frame}
 \frametitle{Source code analyzers}
 \begin{itemize}
 \item Focused on a programming language:
  \begin{itemize}
    \item PyMetrics: \url{http://sourceforge.net/projects/pymetrics/}
    \item This tool provides metrics for Python: Cyclomatic complexity and others
  \end{itemize}
 \end{itemize}
\end{frame}


\begin{frame}
 \frametitle{Source code analyzers}
 \begin{itemize}
 \item Focused on a programming language:
  \begin{itemize}
    \item CCCC: \url{http://cccc.sourceforge.net/}
    \item This tool analyzes the source code of C and C++ providing several metrics
  \end{itemize}
 \end{itemize}
\end{frame}


\begin{frame}
 \frametitle{Source code analyzers}
 \begin{itemize}
 \item Focused on a programming language:
  \begin{itemize}
    \item PerlMetrics: \url{http://www.dunhackin.org/~eisen/perl-metrics/}
    \item This tool analyzes the source code of Perl
  \end{itemize}
 \end{itemize}
\end{frame}




\section{References}

\begin{frame}
 \frametitle{References}
 \begin{itemize}
  \item Study of Available Tools. FLOSSMetrics deliverable 1.1.
 \end{itemize}

\end{frame}

\end{document}